\appendix

\section{Proof Details}\label{sec:proof-details}

\subsection{Coordinate maps, product gauge, and the same-target identity}
\label{app:step-001}

\begin{lemma}[Conditioned left coordinate maps]
\label{lem:step-001-base-coordinate}
Under Assumption~\ref{assump:base_conditioning}, every
$\bar M\in\{\bar A,\bar B,\bar C\}$ has full column rank, and its
coordinate map
\[
\Lambda_M=(\bar M^{\mathsf T}\bar M)^{-1}\bar M^{\mathsf T}
\]
satisfies
\[
\sigma_{\min}(\bar M)\geq\kappa^{-2},
\qquad
\Lambda_M\bar M=I_r,
\qquad
\|\Lambda_M\|_{\mathrm{op}}
=\sigma_{\min}(\bar M)^{-1}\leq\kappa^2.
\]
\end{lemma}

\begin{proof}
The factorization $\bar M=\widetilde M D_{\bar M}$ and
Assumption~\ref{assump:base_conditioning} give
\[
\sigma_{\min}(\widetilde M)\geq\kappa^{-1},
\qquad
\sigma_{\min}(D_{\bar M})
=\min_{j\in[r]}\|\bar m_j\|_2\geq\kappa^{-1}.
\]
Thus, for every $v\in\mathbb R^r$,
\[
\begin{aligned}
\|\bar Mv\|_2
&=\|\widetilde M D_{\bar M}v\|_2\\
&\geq\sigma_{\min}(\widetilde M)\|D_{\bar M}v\|_2\\
&\geq\kappa^{-2}\|v\|_2.
\end{aligned}
\]
Consequently $\sigma_{\min}(\bar M)\geq\kappa^{-2}>0$, so $\bar M$
has full column rank and $\bar M^{\mathsf T}\bar M$ is invertible.  Direct
multiplication then gives $\Lambda_M\bar M=I_r$.

For completeness, if
$\bar M=U\Sigma V^{\mathsf T}$ is an economy singular-value
decomposition, then $U^{\mathsf T}U=V^{\mathsf T}V=I_r$ and $\Sigma$ is
positive diagonal.  Hence
\[
(\bar M^{\mathsf T}\bar M)^{-1}\bar M^{\mathsf T}
=V\Sigma^{-1}U^{\mathsf T},
\]
so its induced Euclidean operator norm is
$\sigma_{\min}(\bar M)^{-1}$.  The preceding lower bound therefore yields
$\|\Lambda_M\|_{\mathrm{op}}\leq\kappa^2$.  This calculation remains
valid at the non-strict conditioning boundary
$\sigma_{\min}(\widetilde M)=\min_j\|\bar m_j\|_2=\kappa^{-1}$.
\end{proof}

\begin{lemma}[Modewise tensor operator in Frobenius geometry]
\label{lem:step-001-tensor-map}
Suppose Assumption~\ref{assump:base_conditioning} holds.

For the setting-ordered map
\[
Q=\Lambda_A\otimes\Lambda_B\otimes\Lambda_C
\]
Lemma~\ref{lem:step-001-base-coordinate} implies
\[
Q(x\otimes y\otimes z)
=(\Lambda_Ax)\otimes(\Lambda_By)\otimes(\Lambda_Cz)
\]
for every $x,y,z\in\mathbb R^n$, and
\[
0<\|Q\|_{\mathrm{op}}
=\|\Lambda_A\|_{\mathrm{op}}
 \|\Lambda_B\|_{\mathrm{op}}
 \|\Lambda_C\|_{\mathrm{op}}
\leq\kappa^6.
\]
Consequently every $R\in\mathbb R^{n\times n\times n}$ obeys
\[
\|QR\|_F\leq\kappa^6\|R\|_F.
\]
\end{lemma}

\begin{proof}
The pure-tensor formula is the defining action of the tensor product in the
fixed mode order.  To compute its norm, take full singular-value
decompositions
\[
\Lambda_M=U_M\Sigma_MV_M^{\mathsf T},
\qquad M\in\{A,B,C\},
\]
where the outer matrices are square orthogonal and each $\Sigma_M$ is
rectangular diagonal.  Under the canonical identification of a third-order
array with the Hilbert tensor product of its mode spaces,
\[
Q=(U_A\otimes U_B\otimes U_C)
  (\Sigma_A\otimes\Sigma_B\otimes\Sigma_C)
  (V_A^{\mathsf T}\otimes V_B^{\mathsf T}\otimes V_C^{\mathsf T}).
\]
The outer maps are Frobenius isometries.  Indeed, on pure tensors this follows
from
\[
\|u\otimes v\otimes w\|_F=\|u\|_2\|v\|_2\|w\|_2,
\]
and expansion in orthonormal product bases extends the identity to the whole
finite-dimensional tensor space.  The middle map is diagonal in product
singular-vector bases, with diagonal entries
\[
\sigma_a(\Lambda_A)\sigma_b(\Lambda_B)\sigma_c(\Lambda_C).
\]
Its largest singular value is therefore
\[
\max_{a,b,c}
\sigma_a(\Lambda_A)\sigma_b(\Lambda_B)\sigma_c(\Lambda_C)
=\prod_{M\in\{A,B,C\}}\|\Lambda_M\|_{\mathrm{op}}.
\]
Every $\Lambda_M$ has rank $r\geq1$ by
Lemma~\ref{lem:step-001-base-coordinate}, so this product is positive.
The same lemma bounds each factor by $\kappa^2$, giving
\[
0<\|Q\|_{\mathrm{op}}\leq(\kappa^2)^3=\kappa^6.
\]
The inequality for an arbitrary $R$ is the definition of the induced
operator norm.  Notice that $Q$ need not be injective on the ambient tensor
space: a component of $R$ may lie in $\ker Q$.  The upper operator inequality
remains valid and is exactly the direction used later.
\end{proof}

\begin{lemma}[Gauge invariance of ambient and coefficient components]
\label{lem:step-001-gauge}
Under Assumption~\ref{assump:base_conditioning} and
Lemma~\ref{lem:step-001-tensor-map}, let $(x,y,z)\in(\mathbb R^n)^3$, put
\[
p=(\Lambda_Ax)\otimes(\Lambda_By)\otimes(\Lambda_Cz),
\]
and let $(x^+,y^+,z^+)$ be the result of the setting-defined componentwise
gauge.  Then
\[
x^+\otimes y^+\otimes z^+=x\otimes y\otimes z
\]
and
\[
(\Lambda_Ax^+)\otimes(\Lambda_By^+)\otimes(\Lambda_Cz^+)=p.
\]
Both statements hold in the positive-norm branch and in every zero-factor
branch.
\end{lemma}

\begin{proof}
Let $u=\|x\|_2$, $v=\|y\|_2$, and $w=\|z\|_2$.  If all three are
positive, set $g=(uvw)^{1/3}$ and
\[
s_x=g/u,
\qquad s_y=g/v,
\qquad s_z=g/w.
\]
The gauge gives $(x^+,y^+,z^+)=(s_xx,s_yy,s_zz)$, and
\[
s_xs_ys_z=\frac{g^3}{uvw}=1.
\]
Multilinearity yields
\[
x^+\otimes y^+\otimes z^+
=(s_xs_ys_z)(x\otimes y\otimes z)=x\otimes y\otimes z.
\]
Linearity of the three coordinate maps gives the same calculation in
coefficient space:
\[
(\Lambda_Ax^+)\otimes(\Lambda_By^+)\otimes(\Lambda_Cz^+)
=(s_xs_ys_z)p=p.
\]

If $uvw=0$, at least one of $x,y,z$ is the zero vector.  The original
ambient rank-one tensor is therefore zero, and the corresponding coordinate
vector is zero by linearity, so $p=0$.  The gauge replaces the triple by
$(0,0,0)$, and both post-gauge tensors are also zero.  The two cases exhaust
all possible norm triples.  Arbitrarily small positive norms do not create a
new case: individual scale factors may be large, but their product remains
exactly one.
\end{proof}

\begin{proposition}[Exact coordinate and same-target interface]
\label{prop:step-001-same-target}
Under Assumption~\ref{assump:base_conditioning} and
Lemmas~\ref{lem:step-001-base-coordinate}--\ref{lem:step-001-gauge}, the
map $Q$ satisfies $0<\|Q\|_{\mathrm{op}}\leq\kappa^6$, the product gauge
preserves every ambient and coefficient component tensor, and every integer
$t\geq0$ satisfies
\[
Q\widehat T_t=C_t\in\mathcal S_t,
\qquad
Q(T-\widehat T_t)=D_r+E_\rho-C_t.
\]
These conclusions include the raw initialization, zero components, the zero
coefficient span, linearly dependent features, and changes in
$\dim(\mathcal S_t)$.
\end{proposition}

\begin{proof}
The operator bound and gauge assertions are
Lemmas~\ref{lem:step-001-tensor-map} and \ref{lem:step-001-gauge}.  For each
$i\in[k]$, Lemma~\ref{lem:step-001-tensor-map} gives
\[
Q(x_{i,t}\otimes y_{i,t}\otimes z_{i,t})
=(\Lambda_Ax_{i,t})\otimes(\Lambda_By_{i,t})
 \otimes(\Lambda_Cz_{i,t})
=p_{i,t}.
\]
Linearity over the finite CP sum therefore gives
\[
Q\widehat T_t
=\sum_{i=1}^kQ(x_{i,t}\otimes y_{i,t}\otimes z_{i,t})
=\sum_{i=1}^kp_{i,t}
=C_t.
\]
By definition, this finite sum belongs to
$\mathcal S_t=\operatorname{span}\{p_{i,t}:i\in[k]\}$.  The conclusion
remains true if some features vanish or are linearly dependent.  If they all
vanish, then $C_t=0\in\mathcal S_t=\{0\}$.

The definition $E_\rho=QT-D_r$ is equivalent to $QT=D_r+E_\rho$.
Subtracting the identity for $Q\widehat T_t$ gives
\[
Q(T-\widehat T_t)=QT-Q\widehat T_t=D_r+E_\rho-C_t.
\]
This is the coefficient image of the same realized ambient residual; no
projected, whitened, base-only, or fixed-subspace surrogate has replaced it.
At $t=0$ the calculation applies directly to the ungauged Gaussian factors.
At later times it applies to the stored factors, while
Lemma~\ref{lem:step-001-gauge} shows that replacing a raw relaxed component
by its gauged representative changes neither tensor sum.  The calculation is
pointwise in $t$ and hence requires neither continuity nor constant rank of
$\mathcal S_t$.

Thus the three modewise $\kappa^2$ bounds combine with no hidden constant,
\[
\|Q\|_{\mathrm{op}}
=\prod_{M\in\{A,B,C\}}\|\Lambda_M\|_{\mathrm{op}}
\leq\kappa^6,
\]
and the exact same-target identity holds at every time.  In the specialization
$Q=I$, $E_\rho=0$, and $T=D_r$, it reduces without loss to
\[
T-\widehat T_t=D_r-C_t,
\qquad C_t\in\mathcal S_t.
\]
All statements in this subsection are deterministic and pointwise.  No
certificate clause or probability conversion is used.
\end{proof}

\subsection{Transport of the coefficient deficit along the projector path}
\label{app:step-002}

\begin{lemma}[One-step transport under projector motion]
\label{lem:step-002-one-step-transport}
For the coefficient Frobenius space and the setting-defined orthogonal
projectors $P_t=\operatorname{Proj}_{\mathcal S_t}$, every integer $t\geq0$
satisfies
\[
\operatorname{dist}_F(D_r,\mathcal S_{t+1})
\geq
\operatorname{dist}_F(D_r,\mathcal S_t)
-\|P_{t+1}-P_t\|_{\mathrm{op}}\|D_r\|_F.
\]
No certificate clause is required for this one-step statement, and it remains
valid for zero, full, stationary, or rank-changing subspaces.
\end{lemma}

\begin{proof}
Let $\mathcal S$ be a coefficient subspace and $P$ its orthogonal projector.
For every $y\in\mathcal S$, the vectors
$(I-P)D_r\in\mathcal S^\perp$ and $PD_r-y\in\mathcal S$ are orthogonal.
Pythagoras therefore gives
\[
\|D_r-y\|_F^2
=\|(I-P)D_r\|_F^2+\|PD_r-y\|_F^2
\geq\|(I-P)D_r\|_F^2.
\]
Equality is attained at $y=PD_r$, so
\[
\operatorname{dist}_F(D_r,\mathcal S)=\|(I-P)D_r\|_F.
\]

Apply this identity at times $t$ and $t+1$.  The exact relation
\[
(I-P_t)D_r=(I-P_{t+1})D_r+(P_{t+1}-P_t)D_r
\]
and the triangle inequality imply
\[
\begin{aligned}
\operatorname{dist}_F(D_r,\mathcal S_t)
&=\|(I-P_t)D_r\|_F\\
&\leq\|(I-P_{t+1})D_r\|_F
   +\|(P_{t+1}-P_t)D_r\|_F\\
&\leq\operatorname{dist}_F(D_r,\mathcal S_{t+1})
   +\|P_{t+1}-P_t\|_{\mathrm{op}}\|D_r\|_F.
\end{aligned}
\]
Rearrangement gives the claimed recurrence.  If $\mathcal S_j=\{0\}$,
then $P_j=0$ and the distance is $\|D_r\|_F$; if $\mathcal S_j$ is the
full coefficient space, then $P_j=I$ and the distance is zero.  A stationary
projector contributes zero charge, while any rank change is charged by its
actual projector difference.  No rank-continuity premise is used.
\end{proof}

\begin{proposition}[Finite-path preservation of the coefficient deficit]
\label{prop:step-002-all-time-deficit}
Suppose clauses 1 and 2 of
$\mathsf C_2(\delta,L_P,\zeta,C_T)$ hold, with
$\delta>0$, $L_P>0$, and $L_P<\delta/4$.  Then, for every integer
$t\geq0$,
\[
\operatorname{dist}_F(D_r,\mathcal S_t)
\geq
\Delta_0-\|D_r\|_F
\sum_{s=0}^{t-1}\|P_{s+1}-P_s\|_{\mathrm{op}}
\geq(\delta-L_P)\|D_r\|_F.
\]
For $t=0$ the sum is empty, and
\[
\delta-L_P>\frac{3\delta}{4}>0.
\]
\end{proposition}

\begin{proof}
We first prove the finite-horizon inequality by induction.  At $t=0$, the
empty-sum convention gives
\[
\operatorname{dist}_F(D_r,\mathcal S_0)=\Delta_0,
\]
so equality holds.  If the claim holds at time $t$, then
Lemma~\ref{lem:step-002-one-step-transport} gives
\[
\begin{aligned}
\operatorname{dist}_F(D_r,\mathcal S_{t+1})
&\geq\operatorname{dist}_F(D_r,\mathcal S_t)
 -\|P_{t+1}-P_t\|_{\mathrm{op}}\|D_r\|_F\\
&\geq\Delta_0-\|D_r\|_F
 \sum_{s=0}^{t}\|P_{s+1}-P_s\|_{\mathrm{op}}.
\end{aligned}
\]
This closes the induction for every finite $t$.

Each projector-motion term is nonnegative.  Clause 2 therefore bounds every
finite prefix by the full path budget:
\[
\sum_{s=0}^{t-1}\|P_{s+1}-P_s\|_{\mathrm{op}}
\leq
\sum_{s=0}^{\infty}\|P_{s+1}-P_s\|_{\mathrm{op}}
\leq L_P.
\]
Clause 1 then yields
\[
\begin{aligned}
\operatorname{dist}_F(D_r,\mathcal S_t)
&\geq\Delta_0-L_P\|D_r\|_F\\
&\geq(\delta-L_P)\|D_r\|_F.
\end{aligned}
\]
No limit of $P_t$, $\mathcal S_t$, or the distance sequence is taken; the
infinite series is used only as a finite-prefix budget.  Finally,
\[
\delta-L_P>\delta-\delta/4=3\delta/4>0.
\]
Since $r\geq1$, $\|D_r\|_F=\sqrt r>0$.  Thus a full coefficient span at
any time, which would have zero distance, is ruled out as a conclusion of the
two clauses rather than excluded as an assumption.  Zero spans, stationary
spans, and rank jumps remain covered by the charged recurrence.  The
projector defect may have the adversarial sign for a lower bound, but its
total all-time charge is exactly at most $L_P\|D_r\|_F$; no cancellation or
decay is invoked.
The implication is deterministic on clauses 1 and 2.  It does not assert
that these clauses, or the full certificate, occur with positive probability.
\end{proof}

\subsection{Transfer to an all-time ambient residual floor}
\label{app:step-003}

\begin{lemma}[Distance to a subspace is one-Lipschitz]
\label{lem:step-003-distance-lipschitz}
In the coefficient Frobenius space, let $\mathcal S$ be any linear subspace
and let $U,V$ be coefficient tensors.  Then
\[
\left|
\operatorname{dist}_F(U,\mathcal S)
-\operatorname{dist}_F(V,\mathcal S)
\right|
\leq\|U-V\|_F.
\]
In particular,
\[
\operatorname{dist}_F(U,\mathcal S)
\geq\operatorname{dist}_F(V,\mathcal S)-\|U-V\|_F.
\]
\end{lemma}

\begin{proof}
For every $S\in\mathcal S$, the triangle inequality gives
\[
\|U-S\|_F\leq\|U-V\|_F+\|V-S\|_F.
\]
Taking the infimum over $S\in\mathcal S$ yields
\[
\operatorname{dist}_F(U,\mathcal S)
\leq\|U-V\|_F+\operatorname{dist}_F(V,\mathcal S).
\]
Interchanging $U$ and $V$ gives the reverse one-sided comparison, and the
two inequalities prove the claim.  The argument includes the zero subspace,
the full coefficient space, and every rank-deficient intermediate subspace.
The constant one is sharp when the perturbation points directly toward or
away from the subspace.
\end{proof}

\begin{proposition}[Coefficient residual floor with explicit smoothing loss]
\label{prop:step-003-coefficient-floor}
Under Assumption~\ref{assump:base_conditioning},
Propositions~\ref{prop:step-001-same-target} and
\ref{prop:step-002-all-time-deficit}, and the clause-4 bound
\[
\|E_\rho\|_F\leq\zeta\|D_r\|_F,
\]
suppose also that $L_P<\delta/4$ and $\zeta<\delta/4$.  Then every
integer $t\geq0$ obeys
\[
\|Q(T-\widehat T_t)\|_F
\geq(\delta-L_P-\zeta)\|D_r\|_F,
\]
and
\[
\delta-L_P-\zeta>\frac{\delta}{2}>0.
\]
\end{proposition}

\begin{proof}
Fix $t\geq0$.  Proposition~\ref{prop:step-001-same-target} gives
\[
Q(T-\widehat T_t)=D_r+E_\rho-C_t,
\qquad C_t\in\mathcal S_t.
\]
Because $C_t$ is an admissible point in $\mathcal S_t$,
\[
\begin{aligned}
\|Q(T-\widehat T_t)\|_F
&=\|D_r+E_\rho-C_t\|_F\\
&\geq\operatorname{dist}_F(D_r+E_\rho,\mathcal S_t).
\end{aligned}
\]
Lemma~\ref{lem:step-003-distance-lipschitz}, with
$U=D_r+E_\rho$ and $V=D_r$, gives
\[
\operatorname{dist}_F(D_r+E_\rho,\mathcal S_t)
\geq
\operatorname{dist}_F(D_r,\mathcal S_t)-\|E_\rho\|_F.
\]
Proposition~\ref{prop:step-002-all-time-deficit} and the displayed
clause-4 bound now yield, without dropping either defect,
\[
\begin{aligned}
\|Q(T-\widehat T_t)\|_F
&\geq(\delta-L_P)\|D_r\|_F-\zeta\|D_r\|_F\\
&=(\delta-L_P-\zeta)\|D_r\|_F.
\end{aligned}
\]
The strict restrictions give
\[
\delta-L_P-\zeta
>\delta-\frac{\delta}{4}-\frac{\delta}{4}
=\frac{\delta}{2}>0.
\]
The time was arbitrary.  Thus the floor holds at the raw entry and after
every update.  The smoothing residual is fixed and subtracted once in each
pointwise comparison; it is not repeatedly accumulated along the horizon.
An adversarially oriented $E_\rho$ may consume the full $\zeta$ allowance,
but no more.
\end{proof}

\begin{proposition}[Exact same-target ambient residual floor]
\label{prop:step-003-ambient-floor}
Under Assumption~\ref{assump:base_conditioning},
Lemma~\ref{lem:step-001-tensor-map},
Proposition~\ref{prop:step-003-coefficient-floor}, and the remaining
clause-4 condition
\[
\|T\|_F\leq C_T\|D_r\|_F,
\qquad C_T>0,
\]
every integer $t\geq0$ satisfies
\[
\|T-\widehat T_t\|_F
\geq
\frac{\delta-L_P-\zeta}{\kappa^6C_T}\|T\|_F.
\]
The bound is deterministic, all-time, and has no hidden constant.
\end{proposition}

\begin{proof}
Fix $t\geq0$ and abbreviate $R_t=T-\widehat T_t$ within this proof.
Lemma~\ref{lem:step-001-tensor-map} gives
\[
0<\|Q\|_{\mathrm{op}}\leq\kappa^6,
\qquad
\|QR_t\|_F\leq\|Q\|_{\mathrm{op}}\|R_t\|_F.
\]
Positivity permits division, while the upper bound on the operator norm gives
the lower-bound direction
\[
\|R_t\|_F
\geq\frac{\|QR_t\|_F}{\|Q\|_{\mathrm{op}}}
\geq\kappa^{-6}\|QR_t\|_F.
\]
This inference needs no injectivity of $Q$ and remains valid when the ambient
residual has components in $\ker Q$.  Proposition~\ref{prop:step-003-coefficient-floor}
therefore yields
\[
\|T-\widehat T_t\|_F
\geq\kappa^{-6}(\delta-L_P-\zeta)\|D_r\|_F.
\]
Since $C_T>0$, the target-scale clause is equivalent to
\[
\|D_r\|_F\geq C_T^{-1}\|T\|_F.
\]
Substitution proves the stated bound.  No division by $\|T\|_F$ occurs, so
$T=0$ is included and gives a zero right-hand side.  Every scale factor is
visible: $L_P$ is the projector-path charge, $\zeta$ is the fixed smoothing
loss, $\kappa^6$ is the coordinate distortion, and $C_T$ is the target-scale
conversion.  There is no probability conversion and no time-dependent
constant.
\end{proof}

\begin{proposition}[Exact/noiseless baseline residual floor]
\label{prop:step-003-baseline-floor}
Under Assumption~\ref{assump:base_conditioning},
Propositions~\ref{prop:step-001-same-target} and
\ref{prop:step-002-all-time-deficit}, suppose the coordinate-orthonormal
exact/noiseless specialization
\[
Q=I,
\qquad E_\rho=0,
\qquad T=D_r
\]
holds.  Then every $t\geq0$ satisfies the stronger bound
\[
\|T-\widehat T_t\|_F\geq(\delta-L_P)\|T\|_F.
\]
\end{proposition}

\begin{proof}
Under the specialization,
Proposition~\ref{prop:step-001-same-target} becomes
\[
T-\widehat T_t=D_r-C_t,
\qquad C_t\in\mathcal S_t.
\]
Consequently
\[
\begin{aligned}
\|T-\widehat T_t\|_F
&=\|D_r-C_t\|_F\\
&\geq\operatorname{dist}_F(D_r,\mathcal S_t)\\
&\geq(\delta-L_P)\|D_r\|_F\\
&=(\delta-L_P)\|T\|_F,
\end{aligned}
\]
where Proposition~\ref{prop:step-002-all-time-deficit} gives the second
inequality.  Thus removing smoothing and coordinate distortion strengthens
the same-target floor instead of replacing it by a conservative remainder,
a stopped statement, or a vanishing-defect surrogate.
\end{proof}

\subsection{Finite represented-tensor variation and objective convergence}
\label{app:step-004}

\begin{lemma}[Finite variation gives ambient convergence]
\label{lem:step-004-finite-variation}
If clause 3 of $\mathsf C_2(\delta,L_P,\zeta,C_T)$ holds, then
$(\widehat T_t)_{t\geq0}$ is Cauchy in ambient Frobenius norm.  Hence there
exists $\widehat T_\infty\in\mathbb R^{n\times n\times n}$ such that
\[
\widehat T_t\longrightarrow\widehat T_\infty
\qquad\text{in Frobenius norm}.
\]
\end{lemma}

\begin{proof}
Define the increment norms and their tails by
\[
v_j=\|\widehat T_{j+1}-\widehat T_j\|_F,
\qquad
V_N=\sum_{j=N}^{\infty}v_j.
\]
Clause 3 states $V_0<\infty$.  Since the series is nonnegative, every tail
is finite and $V_N\to0$.  For integers $s>t\geq0$, exact telescoping gives
\[
\widehat T_s-\widehat T_t
=\sum_{j=t}^{s-1}(\widehat T_{j+1}-\widehat T_j),
\]
and hence
\begin{equation}
\|\widehat T_s-\widehat T_t\|_F
\leq\sum_{j=t}^{s-1}v_j
\leq V_t.
\label{eq:finite-variation-tail}
\end{equation}
Given $\varepsilon>0$, choose $N$ with $V_N<\varepsilon$.  The tails are
nonincreasing, so Equation~\eqref{eq:finite-variation-tail} gives
\[
\|\widehat T_s-\widehat T_t\|_F<\varepsilon
\qquad\text{for all }s>t\geq N.
\]
Thus the represented tensors are Cauchy.  The ambient Frobenius space is
isometric to finite-dimensional Euclidean space $\mathbb R^{n^3}$ and is
therefore complete, which produces $\widehat T_\infty$ in the same tensor
space.  Letting $s\to\infty$ in
Equation~\eqref{eq:finite-variation-tail} also gives the useful tail estimate
\[
\|\widehat T_\infty-\widehat T_t\|_F\leq V_t\longrightarrow0.
\]
This derivation controls only represented tensors.  It requires no factor
bound or factor convergence, Gram conditioning, descent inequality,
design-rank persistence, quotient geometry, or Kurdyka--Lojasiewicz
property.  The first increment $v_0$ is included.  A stationary trajectory
has $v_j=0$ for every $j$ and converges immediately; changing factor
representatives or least-squares ranks does not affect the argument.
\end{proof}

\begin{proposition}[Finite objective limit]
\label{prop:step-004-objective-limit}
Under clause 3 of $\mathsf C_2(\delta,L_P,\zeta,C_T)$ and
Lemma~\ref{lem:step-004-finite-variation}, the objective satisfies
\[
\mathcal L(X_t,Y_t,Z_t)
\longrightarrow
\|T-\widehat T_\infty\|_F^2<\infty.
\]
\end{proposition}

\begin{proof}
For fixed ambient tensors $R,S$, the Frobenius polarization identity and
Cauchy--Schwarz imply
\[
\big|\|T-R\|_F^2-\|T-S\|_F^2\big|
\leq\|R-S\|_F\big(\|T-R\|_F+\|T-S\|_F\big).
\]
Indeed, apply
$\|A\|_F^2-\|B\|_F^2=\langle A-B,A+B\rangle$ to
$A=T-R$ and $B=T-S$.  Taking $R=\widehat T_t$ and
$S=\widehat T_\infty$, and then using the triangle inequality, gives
\[
\begin{aligned}
&\left|
\mathcal L(X_t,Y_t,Z_t)-\|T-\widehat T_\infty\|_F^2
\right|\\
&\quad\leq
\|\widehat T_t-\widehat T_\infty\|_F
\big(\|T-\widehat T_t\|_F+\|T-\widehat T_\infty\|_F\big)\\
&\quad\leq
\|\widehat T_t-\widehat T_\infty\|_F
\left(2\|T-\widehat T_\infty\|_F
+\|\widehat T_t-\widehat T_\infty\|_F\right).
\end{aligned}
\]
The first factor tends to zero by
Lemma~\ref{lem:step-004-finite-variation}; the second tends to the finite
number $2\|T-\widehat T_\infty\|_F$.  This proves convergence to the
displayed finite value.  The conclusion remains valid if the limiting
objective is zero, if $T=0$, if factors diverge while their represented
tensor converges, or if the represented path is stationary.  Mere finiteness
of the variation series supplies no uniform numerical rate, and none is
claimed.
This is a deterministic implication on clause 3; it does not assign a
probability to that clause or to the full certificate.
\end{proof}

\subsection{Conditional theorem closure and the probability boundary}
\label{app:step-005}

\begin{lemma}[Uniform lower bounds pass to a finite limit]
\label{lem:step-005-limit-order}
Let $(a_t)_{t\geq0}$ be a real sequence and let $a,b\in\mathbb R$.  If
\[
a_t\longrightarrow a,
\qquad
a_t\geq b\quad\text{for every }t\geq0,
\]
then $a\geq b$.
\end{lemma}

\begin{proof}
If instead $a<b$, set $\eta=(b-a)/2>0$.  Convergence gives an integer $N$
such that $|a_t-a|<\eta$ for every $t\geq N$.  For such $t$,
\[
a_t<a+\eta=\frac{a+b}{2}<b,
\]
contradicting the pointwise lower bound.  Thus $a\geq b$.  The argument
allows equality, constant sequences, and arbitrary real values of $a$ and
$b$; it uses an ordinary finite limit and makes no extended-real or
convergence-mode upgrade.
\end{proof}

\begin{proposition}[Pathwise closure of the conditional loss floor]
\label{prop:step-005-pathwise-closure}
Under Assumptions~\ref{assump:dimension}, \ref{assump:rank_window},
\ref{assump:base_conditioning}, \ref{assump:gaussian_smoothing}, and
\ref{assump:independent_initialization}, suppose a realized trajectory
belongs to $\mathsf C_2(\delta,L_P,\zeta,C_T)$, where
$\delta,L_P,\zeta,C_T>0$, $L_P<\delta/4$, and
$\zeta<\delta/4$.  Define
\[
m=\delta-L_P-\zeta,
\qquad
\epsilon=\left(\frac{m}{\kappa^6C_T}\right)^2.
\]
Then $m>\delta/2>0$, $\epsilon>0$, and
\[
\lim_{t\to\infty}\mathcal L(X_t,Y_t,Z_t)
=\|T-\widehat T_\infty\|_F^2<\infty,
\]
with
\[
\lim_{t\to\infty}\mathcal L(X_t,Y_t,Z_t)
\geq\epsilon\|T\|_F^2.
\]
\end{proposition}

\begin{proof}
The strict parameter restrictions give the complete margin calculation
\begin{equation}
m
=\delta-L_P-\zeta
>\delta-\frac{\delta}{4}-\frac{\delta}{4}
=\frac{\delta}{2}>0.
\label{eq:positive-margin}
\end{equation}
Since $\kappa\geq1$ and $C_T>0$, this also gives $\epsilon>0$.

Membership in the four-clause certificate supplies the hypotheses of
Propositions~\ref{prop:step-002-all-time-deficit} and
\ref{prop:step-003-ambient-floor}.  Therefore, for every integer $t\geq0$,
\begin{equation}
\|T-\widehat T_t\|_F
\geq\frac{m}{\kappa^6C_T}\|T\|_F.
\label{eq:pathwise-norm-floor}
\end{equation}
Both sides of Equation~\eqref{eq:pathwise-norm-floor} are nonnegative: the
left side is a norm, while the right side is nonnegative by
Equation~\eqref{eq:positive-margin}, $\kappa\geq1$, $C_T>0$, and
$\|T\|_F\geq0$.  Squaring only after this sign check and using the exact
objective definition gives
\begin{equation}
\mathcal L(X_t,Y_t,Z_t)
=\|T-\widehat T_t\|_F^2
\geq
\left(\frac{m}{\kappa^6C_T}\right)^2\|T\|_F^2
=\epsilon\|T\|_F^2
\quad\text{for every }t\geq0.
\label{eq:pathwise-objective-floor}
\end{equation}
No term is absorbed or dropped in this equality.

The same certificate membership supplies clause 3.
Lemma~\ref{lem:step-004-finite-variation} and
Proposition~\ref{prop:step-004-objective-limit} give
\[
\mathcal L(X_t,Y_t,Z_t)
\longrightarrow\|T-\widehat T_\infty\|_F^2<\infty.
\]
Apply Lemma~\ref{lem:step-005-limit-order} to this scalar sequence and the
all-time lower bound in Equation~\eqref{eq:pathwise-objective-floor}.  The
limiting lower bound follows.

The primitive dimension, rank-window, smoothing, and initialization
assumptions preserve the exact theorem domain and joint law; they are not
used to prove certificate membership.  If $T=0$, the proof never divides by
$\|T\|_F$ and the lower bound has zero right-hand side.  Stationary paths are
covered by the finite-variation argument.  No Gram conditioning, factor
boundedness or convergence, descent property, fixed span, basin membership,
design-rank persistence, quotient geometry, or unstated feature of the ALS
map enters the closure.
\end{proof}

\begin{proposition}[Exact conditional event inclusion]
\label{prop:step-005-event-inclusion}
Under Assumptions~\ref{assump:dimension}, \ref{assump:rank_window},
\ref{assump:base_conditioning}, \ref{assump:gaussian_smoothing}, and
\ref{assump:independent_initialization}, let the positive certificate
parameters satisfy $L_P<\delta/4$ and $\zeta<\delta/4$, and define
\[
\epsilon=
\left(\frac{\delta-L_P-\zeta}{\kappa^6C_T}\right)^2.
\]
Then, under the joint smoothing-and-initialization law,
\[
\mathsf C_2(\delta,L_P,\zeta,C_T)
\subseteq
\left\{
\begin{array}{l}
\displaystyle\lim_{t\to\infty}\mathcal L(X_t,Y_t,Z_t)
\text{ exists and is finite},\\[2mm]
\displaystyle\lim_{t\to\infty}\mathcal L(X_t,Y_t,Z_t)
\geq\epsilon\|T\|_F^2
\end{array}
\right\}.
\]
The inclusion asserts neither nonemptiness nor positive probability of the
antecedent event.
\end{proposition}

\begin{proof}
Take an arbitrary outcome of the joint experiment that belongs to
$\mathsf C_2(\delta,L_P,\zeta,C_T)$.  Proposition~\ref{prop:step-005-pathwise-closure}
applies to the realized target and trajectory and proves both properties
defining the event on the right.  Thus every element of the left-hand set is
an element of the right-hand set.

This is an outcome-by-outcome implication inside the primitive probability
space.  It does not invoke a conditional probability, divide by
$\mathbb P[\mathsf C_2]$, or infer a lower bound on that probability.  It
therefore remains valid when the certificate is empty or has probability
zero.  The probability mode is exactly set inclusion, not high probability,
expectation, or an unconditional ALS success statement.
\end{proof}

\begin{proposition}[Exact/noiseless limiting-loss baseline]
\label{prop:step-005-baseline-limit}
Suppose the coordinate-orthonormal exact/noiseless specialization
\[
Q=I,
\qquad E_\rho=0,
\qquad T=D_r
\]
holds, $\delta>0$, $L_P<\delta/4$, and the realized trajectory satisfies
the normalized entry-deficit, finite projector-path, and finite represented-
tensor-variation clauses.  Then
\[
\lim_{t\to\infty}\mathcal L(X_t,Y_t,Z_t)
=\|T-\widehat T_\infty\|_F^2<\infty
\]
and
\[
\lim_{t\to\infty}\mathcal L(X_t,Y_t,Z_t)
\geq(\delta-L_P)^2\|T\|_F^2.
\]
\end{proposition}

\begin{proof}
Proposition~\ref{prop:step-003-baseline-floor} gives, for every $t\geq0$,
\[
\|T-\widehat T_t\|_F\geq(\delta-L_P)\|T\|_F.
\]
The restriction $L_P<\delta/4$ gives
$\delta-L_P>3\delta/4>0$, so both sides are nonnegative and may be squared:
\[
\mathcal L(X_t,Y_t,Z_t)
\geq(\delta-L_P)^2\|T\|_F^2
\quad\text{for every }t\geq0.
\]
The finite-variation clause and
Proposition~\ref{prop:step-004-objective-limit} give the displayed finite
objective limit.  Lemma~\ref{lem:step-005-limit-order} passes the pointwise
baseline floor to that same limit.  Hence the zero-smoothing, identity-
coordinate specialization preserves the stronger constant relative floor;
it is not replaced by an absolute remainder, a stopped statement, or a claim
that only defect terms vanish.
\end{proof}

\subsection{Proof of the Main Theorem}\label{app:main-proof}

\begin{proof}[Proof of Theorem~\ref{thm:main}]
Fix $\kappa,q$ and any fixed constants with the dependence and
strict margins stated in Theorem~\ref{thm:main}.  For arbitrary dimensions,
base triple, target, and ALS trajectory in the theorem scope,
Proposition~\ref{prop:step-001-same-target} proves the exact coordinate
identity and the bound $\|Q\|_{\mathrm{op}}\leq\kappa^6$.
Proposition~\ref{prop:step-002-all-time-deficit} converts the first two
certificate clauses into an all-time coefficient deficit.
Proposition~\ref{prop:step-003-ambient-floor} uses the fourth clause to
transfer that deficit to the exact ambient residual with factor
\[
\frac{\delta-L_P-\zeta}{\kappa^6C_T}.
\]
Independently, Proposition~\ref{prop:step-004-objective-limit} uses the third
clause to produce a finite limit of the same objective sequence.
Proposition~\ref{prop:step-005-event-inclusion} combines those conclusions
and gives precisely the event inclusion and the stated value of $\epsilon$.

The displayed constants have no hidden dependence.  The residual inequality
is valid for every time, while the conclusion concerns the asymptotic
objective; the scalar passage between these modes is
Lemma~\ref{lem:step-005-limit-order}.  The argument is deterministic on the
certificate and performs no probability conversion.

Because the implication holds for every positive constant choice satisfying
the two strict margins, the existential formulation is nonempty at the scalar
level.  For example, one may take
\[
r_0=1,
\qquad C_{\mathrm{dim}}=1,
\qquad \delta=1,
\qquad L_P=\zeta=\frac18,
\qquad C_T=1,
\]
and define $\epsilon$ by the theorem formula.  These are constant functions
of $(\kappa,q)$ and hence have the permitted dependence.  This choice only
discharges the existential quantifiers over scalar constants; it does not
assert that the resulting certificate event is nonempty or has positive
probability.

Finally, Proposition~\ref{prop:step-005-baseline-limit} proves the stated
exact/noiseless baseline conclusion with factor $(\delta-L_P)^2$.  This
completes every part of Theorem~\ref{thm:main}.
\end{proof}
